%! Author = Marjan
%! Date = 19/06/2026

% Preamble
\documentclass[11pt]{book}

% Packages
\usepackage{amsmath}    % For math symbols and equations
\usepackage{graphicx}   % For including graphics
\usepackage{geometry}   % For adjusting page layout
\usepackage{fancyhdr}   % For custom headers/footers
\usepackage{hyperref}   % For hyperlinks
\usepackage{listings}   % For code listing
\usepackage{lipsum}     % For placeholder text (to test layout)
\usepackage{helvet}
\usepackage{arydshln}
\usepackage{tikz}        % For creating diagrams and drawing
\usetikzlibrary{
    positioning,
    arrows.meta,
    shapes,
    fit,
    backgrounds,
    calc
}
\usepackage{pgfplots}    % For plotting graphs
\usepackage{wasysym}
\usepackage{float}
\usepackage{coptic}
\usepackage{xcolor}
\usepackage{tweaklist}
\usepackage{tikzexternal}
\usepackage{xcolor}
\usepackage{listings}
\usepackage{arydshln}
\usepackage{arydshln}

% Define YAML language
\lstdefinelanguage{yaml}{
    keywords={true,false,null,y,n},
    keywordstyle=\color{blue}\bfseries,
    basicstyle=\ttfamily\small,
    comment=[l]{\#},
    commentstyle=\color{green!50!black},
    stringstyle=\color{red},
    morestring=[b]",
    morestring=[b]',
    sensitive=false
}


% Terminal
\lstdefinelanguage{terminal}{
    basicstyle=\ttfamily\small,
    backgroundcolor=\color{black!5},
    frame=single,
    columns=fullflexible
}


% AWS icon command
\newcommand{\awsicon}[4]{
    \node (#1) at #2 {
        \begin{tabular}{c}
        \includegraphics[width=1.8cm]{#3} \\
        #4
        \end{tabular}
    };
}

\tikzset{
    awsarrow/.style={
        -{Latex[length=3mm]},
        thick,
        draw=awsorange
    }
}

%\usepackage{amstex} % Helvetica as a sans-serif font alternative
\renewcommand{\rmdefault}{phv} % Set the default font family to Helvetica

% Set page margins (A4 paper)
\geometry{top=1in, bottom=1in, left=1in, right=1in}

% Set up the header
\pagestyle{fancy}
\fancyhf{}
\fancyhead[L]{Kubernetes Notes}
\fancyhead[C]{Your Name}
\fancyhead[R]{\thepage}
% Document
\begin{document}
% Title Page
    \begin{titlepage}
        \centering
        \vspace*{2in}
        \Huge \textbf{Kubernetes Notes}
        \vfill
        \Large Your Name
        \vfill
        \Large Date: %\today
    \end{titlepage}

    \newpage

    \tableofcontents
    \newpage


    \chapter{Notes info and Introduction}

    \paragraph{Purpose of notes}
    There purpose of these notes are to gain an understanding of Docker and have reference material in case I decide to sit the DCA exam.


    \chapter{Core Concepts}


    \section{Kubernetes Architecture}

    \paragraph{Nodes}
    A node is a machine, physical or virtual upon which Kubernetes is installed.
    This worker machine is where containers will be launched by Kubernetes.
    ....was known as minions in the past...
    These terms may be used interchangeably.

    Consider a node failing.....application will go down.....hence we need multiple nodes.

    \paragraph{Cluster}
    A cluster is a group of nodes grouped together. This way when one node failes, your applicaiton is still accessible form the other nodes.
    Having multiple nodes helps in sharing the load as well.

    \paragraph{Master node}
    Things to consider
    - What is responsible for managing the cluster
    - Where is the information about members of the cluster stored?
    - How are nodes monitored?
    - When a node fails, how do you move the workload of the failed node to another worker node?

    This is the purpose of the master.
    The master is another node with Kubernetes installed in it and is configured as a master.
    The master watches over the nodes in the cluster and is responsible for the actual orchestration of containers on the worker nodes.

    \subsection{What is installed on a system when installing Kubernetes?}
    \begin{itemize}
        \item API Server
        \item etcd
        \item Kubelet
        \item Container Runtime
        \item Controller
        \item Scheduler
    \end{itemize}

    \subsubsection{API Server}
    An API server which acts as the frontend for Kubernetes.
    The user management devices, command line interfaces all talk to the API server to interact with the Kubernetes cluster.

    \subsubsection{etcd keystore}
    Etcd is a distributed, reliable key value store used by Kubernetes to store all data used to manage the cluster.
    The way to think of it is, when you have a multiple nodes and masters in your cluster, etcd stores all that information on all the nodes in the cluster in a distributed manner
    Etcd is responsible for implementing locks within the cluster to ensure that there are no conflicts between the masters,

    \subsubsection{The Scheduler}
    is responsible for distributing work or containers across multiple nodes.
    It looks for newly created containers and assigns them to nodes.

    \subsubsection{Controllers}
    The "brain" behind orchestration.
    They are responsible for noticing and responding when nodes, containers or endpoints go down.
    The controllers decide whether to bring up new containers.

    \subsubsection{Container Runtime}
    They are repsonsible for noticing and responding when nodes, containers or endpoints runtime is the underlying software that is used to run containers.
    An example is Docker, but alternatives exist.

    \subsubsection{Kubelet}
    This is the agent that runs on each node in the cluster.
    The agent is responsble for making sure that hte containrs are running on the nodes as expected.

    \subsubsection{}
    There are two types on server, masters and a worker along with a setup of components which make up Kubernetes.

    ?How does one server become a master and the other the slave.

    The worker node, or minion as it is also know, is where the containers are hosted.

    The master has the Kube API server and that is what makes it a master.
    Similarly the worker nodes have the Kubelet agent that is responsible for interacting with master to provide both health information of the worker node and carry out actions requested by the master on the worker nodes.

    All the information gathered is stored in a key value store on the master.
    The key value store is based on the popular etcd framework.

    The master also has the control manager and the scheduler.

    \subsubsection{Kube command line or kubectl or kube control}
    The kube control tool, is used to deploy and manage applications on Kubernetes cluster to get cluster information, to get the status of other nodes in the cluster. And to manage other things.
    The kubectl clister infor command is used to view infromation about the cluster and the kubectl get nodes command is used to list all the nodes as a part of the cluster.


    \section{Docker-vs-ContainerD}
    Initially kubernets was the container runtime for docker only. Support for other container runtimes was added.
    This was done via an interface called Container Runtime Interface or CRI.

    CRI allows any vendor to work as a container runtime for Kubernetes, as long as they adhere to the OCI standards.
    OCI stands for the open container initiative and it consists of an image spec and runtime spec.

    Image spec means the specifications on how an image should be built.
    An image spec defined the specifications on how an image should be built and the runtime spec defined the standardss on how nay container runtime shoudl be developed.

    Keeping the OCI standards in mind. Anyone cna build a container runtime that can be used by anybody to work with Kubernetes
    Note: Docker wasn't built for CRI. To maintain support between kubernetes and Docker, Kubernetes introduced dockershim.
    Docker runs without CRI.

    Docker isn't just a container runtime along. Docker consists of multiple tools that are put together.
    - Docker CLI, the Docker API and the build tools that help in building images.
    There was support for volumes or security. There is a container runtime called run C and that managed a container called container D
    Container D can be used as a runtime on its own separated from Docker.

    ...dockershim and docker support was removed form kubernetes but docker follows the OCI standards now.

    Container D is a part of Docker but is a separate project on its own.
    You can run container D without having to install Docker itself.
    If you don't need dockers other features you could ideally just install container D alone.

    \subsubsection{CLI - ctr}

    \subsubsection{CLI - nerdctl}
    nerdctl provides a docker like CLI for containerd
    nerdctl supports Docker compose
    nerdctl supported newest features in containerd
    - Encrypted container images
    - Lazy pulling
    - P2P image distribution
    - image signing and verifying
    - Namespaces in Kubernetes

    \subsubsection{CLI-crictl}
    crictl provides a CLI for CRI-compatible container runtimes
    installed separately
    used to inspect and debug container runtimes
    not to create container ideally
    works across different runtimes

%    TODO add table - docker vs crictl


    \section{Pods}
%    TODO create diagram
    Containers are encapsulated into a Kubernetes object know as pods.
    A pod is a single instance of an application.
    A pod is the smallest object you can create in Kubernetes.

    Can increase the number of pods to scale the application.
    Can add a new node and add pods to that if the current nodes do not have capacity.

    Pods usually have a 1 to 1 relationship with
    You do not add containers to a single pod to scale your applicaiton.
%    TODO add diagram

    \subsection{Muti-Container Pods}
    You can have multiple containers in the same pod but they will not be the same kind.
    Example: helper container. The two containers in the same pod can communicate with each other over locahost.
    They can easily share the same storage space and the same network space.


    \section{YAML Basics}


    \section{Pods with YAML}
    Kubernetes uses YAML files as inputs for the creation of objects such as pods, replicas and deployments, services.

    \subsection{YAML in Kubernetes}
    \begin{itemize}
        \item apiVersion: This is the version of the Kubernetes API we are using to create the object. Must use the correct API version for the objects we are trying to create. See other options.
        \item kind: This refers to the type object we are trying to create. Options are Pod, Service, ReplicatSet and Deployment.
        \item metadata: This is data about the object like its name, labels, etc.
        \item spec:
    \end{itemize}


    \begin{lstlisting}[language=terminal,caption={Pod commands}]
        kubectl create -f pod-definition.yml
        kubectl get pods
        kubectl describe pod myapp-pod
    \end{lstlisting}

%TODO update this code listing
    \begin{lstlisting}[language=yaml, caption={Example YAML Configuration}]
        database:
            host: localhost
            port: 5432
            username: admin
            password: secret

        features:
            logging: true
            cache: false

        servers:
            - app01.example.com
            - app02.example.com
    \end{lstlisting}

    \subsection{ReplicatSets}

    \subsubsection{Replication Controller}
    The replication controller helps us run multiple instances of a single pod in the Kubernetes cluster allowing for High Availability.
    This can be done with a single pod. The replication controller can help by automatically brinign up a new pid when the existing one failes.
    The replication controller ensures that hte specified number of pods are running at all times.

    Another reason we need replication controllers is to create multiple pods to share the load across them.
%   TODO the replication controller

    The repilcation controller allows us to create mutiple piods to share the load across them.
    The replication controller spands across multiple notes in the cluster.
    It allows for the application to be scaled when the demand increases.

    Note: there are two similar terms. Replications Controller and Replica Set.
    Both have the same purpose, but they are not the same.
    Replication controller is the older technology that is being replaced by the replicat set.

    Replicat set is the new reccomended way to set up replication.
    There are minor differences in the way each works.

    \subsubsection{How to create a replication controller}
    Start by creating a rc-definition.yml YAML file.
    This will have the standard info as any Kubnernetes file. API, version, kind, metadata and spec.
    The API version is specific to what we are creating.
    The kind is set to replication controller
    Add a name under metadata
    Add labels and assign values.
%    TODO add example
    Add the spec to the YAML. The spec defines what is inside the object we are creating.
    Create a template section under the pod template to be used by the replication controller to create replicas.

    How is the pod template defined?
    Use pod definition file metadata.


    \begin{lstlisting}[language=yaml, caption={Replication controller}]
        apiVersion: v1
        kind: ReplicationController
        metadata:
            name: myapp-rc
            labels:
                app: myapp
                type: front-end
        spec:
            template:
                metedata:
                    name: myapp-pod
                    labels:
                        app: myapp
                        type: front-end
            spec:
                container:
                    - name: nginx-container
                      image: nginx
        replicas: 3
    \end{lstlisting}

    \begin{lstlisting}[language=terminal,caption={Pod commands}]
        kubectl create -f rc-definition.yml
        kubectl get replicationcontroller
        kubectl get pods
    \end{lstlisting}

    \begin{lstlisting}[language=yaml, caption={Replication Set}]
        apiVersion: apps/v1
        kind: ReplicationSet
        metadata:
            name: myapp-replicaset
            labels:
                app: myapp
                type: front-end
        spec:
            template:
                metedata:
                    name: myapp-pod
                    labels:
                        app: myapp
                        type: front-end
            spec:
                container:
                    - name: nginx-container
                      image: nginx
        replicas: 3
        selector:
            matchLabels:
                type: front-end
    \end{lstlisting}

    The selector section helps the replica set identify what pods fall under it.
    Replica sets can also manage pods that were no created as a part of the replica set creation.
    Example, pods created before the creation of the replica sets that match the labels specified in the selector.
    The replicaset will also take those pods into consideration when creating the replicas.

    The selector is one of the major differences between replication controller and replica set.
    The selector is not a required field in case of a replication controller, but it is still available.

    Match label selector simply matches the labels specified under it to the labels on the pods.

    \subsubsection{Labels and selectors}
    Uses of replica sets. Use them to monitor existing pods if they have already been created.

    The role of the the replica set is to monitor the piods and if any of them were to fail, deploy new ones.
    The replica set is a process that monitors the pods.

    Labelling the pods is handy since we can now use these as a filter for replica set under the selector section, we use the match labels filter and provide the same filter we used while creating the pods.

    \subsubsection{Scaling with a replica set}
    One way is to update the number of replicas in the definition file to six replicas the run kubectls apply command to sspecifiy the same file using the f parameter and the will update the replica set to have six replicas.

    Another way is to run the kube control scal command. Use the replicas parameter to provide the number of replicas and specify the same file as input.
    - You can either input the definition file or provide the replica set name in the type name format. However, remember that using hte file name as input will not result in the number of replicas being updated automatically in the file.
    --- In other words the number of replicas in the replica set definition fil will still be three, even through via this command it was scaled to siz.

    There also options available for automatically scaling the replica set based on load, but that is an advanced topic.....

%    TODO add command review.
    \begin{lstlisting}[language=terminal,caption={Pod commands}]
        kubectl create -f replicaset-defintion.yml
        kubectl get replicaset
        kubectl delete replicaset myapp-replicaset
        kubectl replace -f replicaset-defintion.yml
        kubectl scale -replicas=6 -f replicaset-definition.yml
    \end{lstlisting}

    \subsubsection{Deployment}
    Each container is encapsulated inside a pod.
    Multiple pods are deployed using replication controllers or replica sets.

%    TODO diagram
    Higher in the hierarchy of Kubernetes, above replica sets is deployments

    \begin{lstlisting}[language=yaml, caption={Replication Set}]
        apiVersion: apps/v1
        kind: Deployment
        metadata:
            name: myapp-replicaset
            labels:
                app: myapp
                type: front-end
        spec:
            template:
                metedata:
                    name: myapp-pod
                    labels:
                        app: myapp
                        type: front-end
            spec:
                container:
                    - name: nginx-container
                      image: nginx
        replicas: 3
        selector:
            matchLabels:
                type: front-end
    \end{lstlisting}

    \begin{lstlisting}[language=terminal,caption={Pod commands}]
        kubectl create -f deployment-defintion.yml
        kubectl get deployments
        kubectl get replicasets
        kubectl get pods
        kubectl get all
    \end{lstlisting}

    \subsubsection{Namespaces}

    \subsubsection{Kubectl Explain Command}
    Pods, deployments and clusters are deployed within a namespace.

    Kubernetes creates a default namespace when the cluster is setup and these are created under a namespace called kubesystem.
    The kube-public namespace is where resource that should be made available to all users are created.

    When a kubernetes cluster grows to production/enterprise level you will want to consider using namespaces.

    You can create your own namespaces as well.
    Example, if you want to use the same cluster for both dev and production environments but at the same time isolate the rousources ebtween them.
    You can create a different namespace for each of them.

    That way while working in the dev environment, you don't accidentally modify resources in production.

    Each of these namespaces can hav its own set of policies that define who can do what.

    You can also assign a quota of resources to each of these namespaces.
    That way, each namespace is guaranteed a certain amount and does not use more that its allowed limit.

    Resources within a default namespace can refer to each other by their names.

    You can make sure that the pod gets created into the development environment all the time by moving the namepsace intot hte pod definitno file.

    \begin{lstlisting}[language=yaml, caption={Replication Set}]
        apiVersion: apps/v1
        kind: Deployment

        metadata:
            name: myapp-pod
            namespace: dev
            labels:
                app: myapp
                type: front-end
        spec:
            containers:
            - name: niginx-container
              image: nginx
    \end{lstlisting}

    \begin{lstlisting}[language=terminal,caption={Pod commands}]
        kubectl get pods
        kubectl get pods --namespace=kube-system

        kubectl create -f pod-definition.yml
        kubectl create -f pod-definition.yml --namespace=dev
    \end{lstlisting}

    \subsubsection{Create Namespace}

    \begin{lstlisting}[language=yaml, caption={Replication Set}]
        apiVersion: apps/v1
        kind: Namespace
        metadata:
            name: dev
    \end{lstlisting}

    \begin{lstlisting}[language=terminal,caption={Pod commands}]
        kubectl create -f namespace-dev.yml
        kubectl create namespace dev
        kubectl get pods --namepsace=dev
        #Used to change context for namespace command
        kubectl config set-context $(kubectl config current-context) --namespace=dev

        kubectl get pods
        kubectl get pods --namespace=default
        kubectl get pods --namespace=prod

        # view pods in all namespaces
        kubectl get pods --all-namespaces
    \end{lstlisting}

    \subsubsection{Resource Quota}
    \begin{lstlisting}[language=yaml, caption={Replication Set}]
        apiVersion: v1
        kind: ResourceQuota
        metadata:
            name: compute-quota
            namespace: dev

        spec:
            hard:
                pods: "10"
                request.cpu: "4"
                requests.memory: 5Gi
                limits.cpu: "10"
                limits.memory: 10Gi
    \end{lstlisting}

    \begin{lstlisting}[language=terminal,caption={Pod commands}]
        kubectl create -f compute-quota.yaml
    \end{lstlisting}

    \subsubsection{Services}
    Kubernetes services enable communication between various components within and outside of the application.
    These services help connect application together with other applicaitons or users.

    Example to consider. Application has a group of pods running variaous sections such as group for running back end processes, a group for frontend services and a gruop connecting to an external data source.

    Services enabled the front end applicaiton to made avaiable to end users to end users.
    It helps communication between backend and fronend pods and helps in establishing connectivity to an external data source.

    Services snable loose coupling between microservices in out application.

    A kubernetes service is an object just like pods, replica sets or deployments that we worked with before.
    One of its use cases is to listen to a port on the node and forward requests on that port to a port on the pod running the web application.

    This type of service is known as a node port services, because the service listens to a port on the node and forwards the request to the pods.

    \subsubsection{Service Types}
    There are different types of services.

    The first type is the node port where the service makes an internal port accessible on a port on the node.

    The second is cluster IP and in this case the service creates a virtual IP inside the cluster to enable communication between different services, such as a set of front end servers and a set of backend servers.

    The third type is a load balancer, where it provisions a load balancer for our applicaiton in supported cloud providers.

    \subsubsection{Node Port}
%    TODO see diagram.
    There are three ports involved.
    The port on the pod where the actual web server is running. This is referred to as the target port because that is where the service forwards the request.
    The second port is the port on the service. Referred to as the port. This service is in fact like a virtual server inside the node, inside the cluster.
    It has its own IP address and that IP address is called the cluster IP of the service.
    Finally, we have the port on the node itself which we can use to access the web server externally.
    -> this is known as the node port.
    Node ports can only be in a valid range which is by default from 30,000 to 32,767.

    \subsubsection{How to create a service}
    Use a definition file to create a service. The high level structure of the file remains the same as before.

    \begin{lstlisting}[language=yaml, caption={Replication Set}]
        apiVersion: v1
        kind: Service
        metadata:
            name: myapp-service

        spec:
            type: NodePort
            ports:
                - targetPort: 80 #Note that this is an array - you can have multiple port mappings in a single service
                  port: 80 #Out of all of these this is the mandatory port
                  nodePort: 30008
            selector: # Use to identify the pod
                app: myapp
                type: front-end
    \end{lstlisting}

    If you don't provide a target port, it is assumed to be the same as port.
    Also, if you don't provide a node port, a free port in the valid range between 30,000 and 32,767 is automatically allocated.

    \begin{lstlisting}[language=terminal,caption={Pod commands}]
        kubectl create -f service-definition.yml
        kubectl get services
        # We can now this port to access the web service using Curl or a web browser
        curl http://192:168.1.2:30008
    \end{lstlisting}

%    TODO multiple pods diagram wihtin a cluster
    Requests are randomly routed to pods....
    The service acts as a built-in load balancer to distribute load across different pods

%    TODO - diagram about when pods are distribted across multiple nodes.
    In this case we have web applications in pods on separate nodes in a cluster.
    Kubenetes creates a service which spans all the nodess in the cluster.

    They can be accessed using curl with the correct IP address and port.

    To summarise in either case the service is created exactly the ame without having to do any additional steps during hte service creation.

    When pods are removed or added, the service is automatically updated, making it highly flexible and adaptive.
    Once created, you won't typically have to make any additional configuration changes.

    \subsubsection{Services - Cluster IP}

    A full stack web application typically has different kinds of pods hosting different parts of an appl,.ication.
    You may have a number of pods running a frontend web server, another set of pods running a backend server, a set of pods running a key value store like reddis etc.

    They will need to communicate with each other.

%     TODO add diagram
    What is the right way for them to communicate?

    The pods all have an IP address assigned to them but these IPS as we know are not static.
    These pods can go down any time and new pods are created all the time.

    You cannot rely on these IP address for internal communication between the application.

    A Kubernetes service can help us group the pods toghether and provide a single interface to acces the pods in the group.

    Example:
    A service created for the backend pods will help group all the backend pods together andd provide a single interface for other pods to access this service.
    The requests are forwarded to one of the pods under ther service randomlyu.

    Similary create additional service for readings and allow the backend pods to access th reading systems through the service.

    This enables us to easily and effectively deploy a microservices based application on Kubernetes cluster.

    Each layer can now scale or move as required without impacting communication between various services.

    Each service gets an IP and name assigned to it inside the cluster and that is the name that should be used by other paths to access the service.

    This type of service is known as the cluster UIIP

    \begin{lstlisting}[language=yaml, caption={Cluster IP service definition}]
        apiVersion: v1
        kind: Service
        metadata:
            name: back-end
        spec:
            type: ClusterIP # This is default
            ports:
              - targetPort: 80 # Where the backend is exposed
                port: 80 # Where the service is exposed
            selector: # To Link a service to a set of ports
                app: myapp
                type: back-end
    \end{lstlisting}

    \begin{lstlisting}[language=terminal,caption={Commands}]
        kubectl create -f service-definition.yml
        kubectl get services
    \end{lstlisting}

    \subsubsection{Kubectl Command}
    SEE DOCS kubectl cheat sheet and kubectl overview

    \begin{lstlisting}[language=terminal,caption={Commands}]
        kubectl api-resources
        kubectl explain pods #top level
        kubectl explain pods --recursive # Easier to create YAML files out of this
    \end{lstlisting}


    \chapter{Configuration}


    \section{Define build and modify container images}
    Why create your own image?
    BEcause you cannot find a component or a service that you want to use as a part of your application on Docker Hub already
    or your team decided that you're development will be dockerized for ease of shipping and deployment.

    \subsubsection{How to create my own image?}
    First we need understand what we are containerizing or what application we are creating and imasge ofr and how the appliicaiton is built.
    Start by thinking if you want to deploy the application manually.

    \begin{itemize}
        \item OS - Ubuntu
        \item Update apt repo
        \item Install dependencies using apt
        \item Install Python dependencies using pip
        \item Copy source code to /opt folder
        \item Run the web server using "flask" command
    \end{itemize}

    \begin{lstlisting}[language=yaml, caption={Example docker file}]
        FROM Ubuntu # Base OS or another image

        RUN api-get update # Install all dependencies
        RUN apt-get install python

        RUN pip install flask
        RUN pip install flask-mysql

        COPY . /opt/source-code # Copy source code to a location

        ENTRYPOINT FLASK_APP=/opt/source-code/app.py flask run # Specify Entrypoint - specify a command that will be run when the image is run as a container
    \end{lstlisting}

    \begin{lstlisting}[language=terminal,caption={Commands}]
        docker build Dockerfile -t username/my-custom-app # The above command will create an image locally on the system.
        docker push username/my-custom-app #This makes the image available on the local docker hub
    \end{lstlisting}

    \subsubsection{Dockerfile}
    A docker file is a text file written in a specifc format.
    <Instruction> <Argument>

    Each of these instruct docker to perform a specific file while creating the image.

    When docker builds the images it builds this in a layered architecture.
    Each line of instruction creates a new layer in the docker image with just the changes from the previous layer.

    When running the "docker build ." command you can see the various steps involved and the result of each task.
    A docker build can be restarted from a step incase a step fails or if you were to add new steps in the build process you wouldn't have to start again.
    Previous layers are reused from a cache.

    This helpful especially when you update source of your application, as it may change more frequently.
    Only layers above the updated layers needs to be rebuilt.

    Everything can be containerized including browsers like firefox etc and utilities like curl.


    \section{Commands and Arguments in Docker}
    note: not required for cert but useful in usage

    Unlike virtual machine, containers are no meant to host an operating system
    Containers are meant to run a specific task or process, such as to host an instance of a web server, or application server or a database.
    Or simply to carry out some kind of computation or analysis.

    Once the task is complete the container exits.
    A container only lives as long as the process inside it is alive.

    If the web service inside the container is stopped or crashses the container exists.

    Command line arguments can be appended to the entry point.


    \section{Commands and Arguments in Kubenetes}

    \begin{lstlisting}[language=yaml, caption={Cluster IP service definition}]
        apiVersion: v1
        kind: Pod
        metadata:
            name: ubuntu-sleeper-pod
        spec:
            containers:
                - name: ubuntu-sleeper
                  image: ubuntu-sleeper
                  command: ["sleep.2.0"] # Corresponds to entry point in docker file
                  args: ["10"] # Corresponds to CMD in docker file
    \end{lstlisting}
%    TODO see note

    \begin{lstlisting}[language=terminal,caption={Commands}]
        docker run -r APP_COLOR=pink simple-webapp-color
    \end{lstlisting}


    \section{Envrinoment Variables}
    \begin{lstlisting}[language=yaml, caption={Example pod definition file with yml}]
        apiVersion: v1
        kind: Pod
        metadata:
            name: simple-webapp-color
        spec:
            container:
            - name: simple-webapp_colo
            image: simple-webapp-color
            ports:
            - containerPort: 8080
            env: # These are plain key values
            - name: APP_COLOR
            value: pink
    \end{lstlisting}

    There are other ways of setting these environment values. Plain key values and secrets

    \begin{lstlisting}[language=yaml, caption={Example pod definition file with yml}]
        env:
            - name: APP_COLOR # Config map
            valueForm:
                configMapKeyRef:

        ////////
        env: # Secrets
            - name: APP_COLOR
            valueFrom:
                secretKeyRef:

    \end{lstlisting}


    \section{ConfigMaps}

    used to manage environment variable centrally

    \begin{lstlisting}[language=yaml, caption={Example pod definition file with yml}]
        apiVersion: v1
        kind: pod
        metadata:
            name: simple-webapp-color
        spec:
            containers:
            - name: simple-webapp-color
            image: simple-webapp-color
            ports:
                - containerPort: 8080
            envFrom:
                - configMapRef:
                    name: app-config
    \end{lstlisting}

    \begin{lstlisting}[language=yaml, caption={Example pod definition file with yml}]
        APP_COLOR: blue
        APP_MODE: prod
    \end{lstlisting}

    Create the ConfigMap
    Inject the confifg map into the pod

    There are two ways of creating a config map.

    The impoerative way without using a config map definition file - kubectl create configmap
    The declarative way by using a config map definition file.

    If you don't wise to create a config map definition file, you just use the Kubectl control command

    \begin{lstlisting}[language=yaml, caption={Example pod definition file with yml}]
        kubectl create configmap <config-name> \
            --from-literal=APP_COLOR=blue \
            --from-literal=APP_MOD=prod
    \end{lstlisting}

    --- using a path to file

    \begin{lstlisting}[language=yaml, caption={Example pod definition file with yml}]
        kubectl create configmap <config-name> \
            --from-file=<path-to-file>
    \end{lstlisting}

    \subsubsection{Using declarative approach}
    \begin{lstlisting}[language=yaml, caption={Example pod definition file with yml}]
        kubectl create -f
    \end{lstlisting}

    \begin{lstlisting}[language=yaml, caption={Example pod definition file with yml}]
        apiVersion: v1
        kind: ConfigMap
        metadata:
            name: app-config
        data:
            APP_COLOR: blue
            APP_MODE: prod
    \end{lstlisting}

    \begin{lstlisting}[language=yaml, caption={Example pod definition file with yml}]
        kubectl create -f config-map.yml
    \end{lstlisting}

    \begin{lstlisting}[language=yaml, caption={Viewing config maps}]
        kubectl get configmaps
        kubectl describe configmaps
    \end{lstlisting}

    \subsubsection{Example}
    \begin{lstlisting}[language=yaml, caption={Example pod definition file with yml}]
        apiVersion: v1
        kind: ConfigMap
        metadata:
            name: app-config
        data:
            APP_COLOR: blue
            APP_MODE: prod
    \end{lstlisting}

    \begin{lstlisting}[language=yaml, caption={Example pod definition file with yml}]
        apiVersion: v1
        kind: Pod
        metadata:
            name: simple-webapp-color
        spec:
            containers:
            - name: simple-webapp-color
            image: simple-webapp-color
            ports:
                - containerPort: 8080
            envFrom:
                - configMapRef:
                    name: app-config
    \end{lstlisting}

    \begin{lstlisting}[language=yaml, caption={Viewing config maps}]
        kubectl create -f pod-definition.yaml
    \end{lstlisting}


    \section{Secrets}
    \begin{lstlisting}[language=yaml, caption={Example pod definition file with yml}]
        import os
        from flask import Flask

        app = Flask(__name__)

        @app.route("/")
        def main():
            mysql.connector.connect(host='mysql', database='mysql', user='root', password='paswrd')
            return render_template('hello.html', color=fetchcolor())

        if __name__ = "__main__":
            app.run(host="0.0.0.0", port="8080")
    \end{lstlisting}

    \begin{lstlisting}[language=yaml, caption={Example pod definition file with yml}]
        apiVersion:
        kind: ConfigMap
        metadata:
            name: app-config
        data:
            DB_Host: mysql
            DB_User: root
            DB_Password: paswrd
    \end{lstlisting}

    Note the hardcoded values. This is bad.
    Can store this in a Config.
    Use secrets to store this sensitive information.
    Similar to config maps except that they are stored in an encoded or hashed format

    Two steps when working wtih secrets
    - Create the secrets
    - Inject into the pod

    \begin{lstlisting}[language=yaml, caption={Example pod definition file with yml}]
        DB_Host: mysql
        DB_User: root
        DB_Password: paswrd
    \end{lstlisting}

    \begin{lstlisting}[language=yaml, caption={Imperative creation}]
        kubectl create secret generic
    \end{lstlisting}

    \begin{lstlisting}[language=yaml, caption={Declarative creation}]
        kubectl create -f
    \end{lstlisting}

%    TODO see example of imperative command

%    TODO see example of Declarative command

    \begin{lstlisting}[language=yaml, caption={Declarative creation}]
        kubectl create -f
    \end{lstlisting}

    \begin{lstlisting}[language=yaml, caption={Example pod definition file with yml}]
        apiVersion:
        kind: Secret
        metadata:
            name: app-secret
        data:
            DB_Host: mysql #This is plain text. Data must be specified in an encoded form
            DB_User: root
            DB_Password: paswrd
    \end{lstlisting}

    \begin{lstlisting}[language=yaml, caption={Declarative creation}]
        kubectl create -f secret-data.yaml
    \end{lstlisting}

    \subsubsection{Encoding Secrets}

    \begin{lstlisting}[language=yaml, caption={Example pod definition file with yml}]
            DB_Host: mysql #This is plain text. Data must be specified in an encoded form
            DB_User: root
            DB_Password: paswrd
    \end{lstlisting}

    \begin{lstlisting}[language=yaml, caption={Declarative creation}]
        echo -n 'mysql' | base64
        echo -n 'root' | base64
        echo -n 'paswrd' | base64

        kubectl get secrets # view secrets
        kubectl describe secrets
        kubectl get secret app-secret -o yaml

        echo -n 'mysql' | base64 --decode
        echo -n 'root' | base64 --decode
        echo -n 'paswrd' | base64 --decode
    \end{lstlisting}

    \begin{lstlisting}[language=yaml, caption={pod-definition.yml}]
        apiVersion: v1
        kind: pod
        metadata:
            name: simple-webapp-color
        labels:
            name: simple-webapp-color
        spec:
            containers:
            -name: simple-webapp-color
            image: simple-webapp-color
            ports:
            - containerPort: 8080
            envFrom:
            - secretRef:
                name: app-secret
    \end{lstlisting}

    \begin{lstlisting}[language=yaml, caption={secret-data.yaml}]
        apiVersion: v1
        kind: Secret
        metadata:
            name: app-secret
        data:
            DB_host: ewafe-
            DB_User: ;afweohiu
            DB_Password: ;oiuhj
    \end{lstlisting}

    \begin{lstlisting}[language=yaml, caption={Declarative creation}]
        kubectl create -f pod-definition.yaml
    \end{lstlisting}

    There are other ways to inject secrets into pods.
    You can inject single environment variables or inject the whole secret as files in a volume.

    If you were to mount the secret as a volume in the pod, each attribute in the secret is created as a file with the value of the secret as its content.

    \begin{lstlisting}[language=yaml, caption={secret-data.yaml}]
        volumes:
        - name: app-secret-volumes
        secret:
            secretName: app-secret
    \end{lstlisting}

    \begin{lstlisting}[language=yaml, caption={Commands to run from within the container}]
        ls /opt/app-secret-volumes
        cat /opt/app-secret-volumes/DB_Password
    \end{lstlisting}

    If you were to mount the secret as a volume in the pod, each attributed in the secret is created a separate file with the value of the secret as its content.


    \section{Pre-requesite for Kubernetes - Security in Docker}
    Consider a host with Docker installed on it.
    This host has a set of its own processes running such as a number of operating system processes, the docker daemon itself, the SSH server, etc.
    %See diagram

    Unlike virtual machines containers are no completely isolated from their host containers and the hosts share the same kernel.

    Containers are isolated using namespaces.

    In Linux, the host has a namespace and the containers hvae their own namepsace.
    All the processes run by the containers are in fact run on the host itself, but in their own namespace.

    As fare as the Docker container is concerned, it is in its own namespace and it can see its own processes only.
    --> If cannot see anything outside of it or in any other namespace.

    For the Docker host, all processes of its own, as well as those in the child namespaces are visiable as just another process in the system.

    When you list the procfefss on the host, you see a list of processes.
    ---> Profcsess can have different process IDs in different namespaces.
    ---> Thats how Docker isolates containers within a system.

    That is process isolation.

    \subsubsection{Users in the context of security}
    The docker host has a set of users, a root user as well as a number of non-root users.
    By default, Docker runs processes within containers as the roots user.
    The prrocess is run as the root users.
    This can be seen via commands both inside the container and outside the container.

    If you do not want the process within the container to run as the root user, you may set the users using the user option within the docker run command and specify the new user ID
    ---- The process now runs with the new user ID.

    Another way to enforce user security is to have thiss defined in the Docker image itselfl at the time of creation.

    \subsubsection{Root user}
    Docker implements a set of security features that limits the abilities of the root user within the container.
    So the root user within the container isn't really like the root user on the host.
    Docker uses Linux capabilities to implement this.
    Note: Root users can literally do anything, as can a process run by the root users.
    It has unrestricted access to the system from modifying files and permissions on files, access controller, creating or killing processes, setting group ID or user ID
    --> Performing networrk related operations such as binding to network ports, broadcasting on a network, controlling network port system related operations like rebooting the host, manipulating the system clock and others.

    All of these are the different capabilities on a Linux system and you can see a full list of this at this location.
    /usr/include/linux/capability.h

    You can contorl and limit what capabilities are made avaiable to a user.
    By default, Docker runs a container with a limited set of capabilities and so the processes running within the container do not have the privilleges to reboot the host or perform other operations that can disruipt the host or other containers runnign on the same host.

    If you wish to voerride this behaviour and provide additional privilges that what is avaialbe you can use

    \begin{lstlisting}[language=terminal,caption={Commands}]
        docker run --cap-add MAC_ADMIN ubuntu
        docker run --cap-drop KILL # drop privilleges
        docker run --privileged #Run a container with all privilleges enabled
    \end{lstlisting}


    \section{Security Contexts}
    In Kubenerntes containers are encapsulated in pods. You may choose to confgiure the securtiy settings at a container level or at a pod level.
    If you configure it as both the pod and the container level the settings on the container will override the settings on the pod.

    \begin{lstlisting}[language=yaml, caption={secret-data.yaml}]
        apiVersion: v1
        kind: Pod
        metadata:
            name: web-pod
        spec:
            securityContext:
                runAsUser: 1000 # Pod level user id
            containers:
                - name: ubuntu
                image: ubuntu
                command: ["woop", "3768"]
                securityContext:
                    runAsUser: 1000 # Container level user id security
                    capabilities:
                        add: ["MAX_ADMIN"] #Specify a list of capabilities to add to the pod. ---- Note capabilities are only suported at the container level and not at the POD level.
    \end{lstlisting}


    \section{Resource Requirements}
%TODO see diagram
    Each node has a set of CPU and memory resource available.
    Every pod requried a set of resources to run.
    Whenever a pod is placed on a node, it consumes the resources available on that node.
    It is the Kubernetes Scheduler that decided which node a pod goes to.
    The scheduler takes into consideration the amount of resources requried by a pod and those available on the nodes.
    It identifies the best node to place a pod on.

    If node have no sufficicent resources available. The scehduler avoids placing the pod on those nodes.
    Instead places the pod on the one where sufficient resources are available.
    If there are no sufficient resource availalbe on any of the nodes then the scehduler holds back, scheduling the pod and you will see the pod in a pending state.

    If you look at the events using the kube control describe pod command you will see that there is an insufficient CPU command.

    You can specify the amoutn of CPU and memeory required for a pod when creating one.
    This is known as the resource request for a container.

    When the scheduler tries to place the pod on a node, it uses these number to identify a node which has sufficent amount of reosources availalbe.

    \subsubsection{Example Yaml with resource request}
    \begin{lstlisting}[language=yaml, caption={pod-definition.yaml}]
        apiVersion: v1
        kind: Pod
        metadata:
            name: simple-webapp-color
            labels:
                name: simple-webapp-color
        spec:
            containers:
            - name: simple-webapp-color
              image: simple-webapp-color
              ports:
                - containerPort: 8080
              resources:
                requests:
                    memory: "4gi"
                    cpu: 2
    \end{lstlisting}

    When the scheduler gets a request to place this pod, it looks for a node that has this amount of resources available.
    And when the pod gets placed on a node, the pod gets a guaranteed amount of resources available for \textbf{it}

    What does one count of CPU really mean?
    0.1 can be expressed as 100M where m standard for Milli
    you can go as low as 1milli

    1 count of CPU is equivalent to 1 AWS vCPU , 1 GCP core, 1 Azure Core, 1 Hyperthread <----

    Can request a higher levle of CPU provded that youru nodes are sufficiently funded.

    Similarly you can specify the memeory using the Mi siffix ro specify the memeory using the exact amount....

    1G - base 10
    1Gi - base 2

    A container has no limit to the amount of resources it can consumer on a node.
    It can start with on CPU and consumer as many resources as it requires and that can suffocate the native process on the node or other containerss or resources.
    You can set a limit of the resource usage on these pods.
    You can limit it to 1 vCPU or limit it to 512 Mi

    \begin{lstlisting}[language=yaml, caption={pod-definition.yaml}]
        apiVersion: v1
        kind: Pod
        metadata:
            name: simple-webapp-color
            labels:
                name: simple-webapp-color
        spec:
            containers:
            - name: simple-webapp-color
              image: simple-webapp-color
              ports:
                - containerPort: 8080
              resources:
                requests:
                    memory: "4gi"
                    cpu: 2
                limits: # Resources limit
                    memory: "2Gi"
                    cpu: 2
    \end{lstlisting}

    Now when the pod is created, Kubernetes sets new limits for the container.
    Rememeber that the limits and reqeusts are set for each container within a pod.
    If there are multiple containers, then each container can have a request or limited set of its own.

    What happends when a pod tries to exceed the reousrces beyond its specified limit?
    In case of the CPU, the system throttles the CPU so that it does not go beyone the specified limit.
    A container cannot use more CPU resources that its limit.

    --> This is not hte case with memeory.
    A container can use more memory reousrces that its limit, so it if a pod tries to consumer moremorry that its limit constanfly.
    --> The pod will be terminated and you will seet that with the pod terminated command with an error in the logs.
    --> can also be ssen in the outpout of the describe command when you run it

    OOM - referes to an out of memory container kill.

    By default kubernetest does not have a CPU or a memeory request or limit set.
    This means that any pod can consume as much resource s as required on any node and suffocate other popds or porecess that are running on the node, of resources.
    This is important to not.

    Examples scenario: %TODO see diagrams
    Say there are two pods competing for CPU reousrces on the cluster. POD/Container.
    Withouth a resource or limit set, one pod can consume all the CPU resources on the node and prevent the second part of the required reousrces.

    Can set limits to prevent users from misusing the infrastructure.

    Need make sure that pods have request set because that is the only way to guarantee resources.

    For memeory it is quite similar. %TODO see diagrams
    Unlike memeory we cannot throttle it. Once memory is assigned the only way to free it is to kill it.

    \subsubsection{Limit Ranges}
    Limit ranges can help you define default values to be set for containers in pods that are created withouth a request or limit specified in the Pod definitino files.
    This is applicable at the namespace level.

    \begin{lstlisting}[language=yaml, caption={Limit range example}]
        apiVersion: v1
        kind: LimitRange
        metadata:
            name: cpu-resource-constraint
        spec:
            limits:
            - default:
                cpu: 500m
              defaultRequest:
                cpu: 500m
              max:
                cpu: "1"
              min:
                cpu: 100m
              type: Container
    \end{lstlisting}

    \begin{lstlisting}[language=yaml, caption={Limit range example}]
        apiVersion: v1
        kind: LimitRange
        metadata:
            name: memory-resource-constraint
        spec:
            limits:
            - default:
                cpu: 1Gi
              defaultRequest:
                cpu: 1Gi
              max:
                cpu: 1Gi
              min:
                cpu: 500Mi
              type: Container
    \end{lstlisting}

    These limits are enforced when a pod is created.
    If you create or change a limit rang it does not affect existing pods.
    It will only affect newer pods that are crated after the limit range is created or updated.

    There is a way to restrict the total amount of resources that can be consumed by application in a Kubernetes cluster.
    We can create quotas at a namespace level.
    So a resource quota is a namepsace laveal object that can be created to set hard limits for requests and limits.

    \begin{lstlisting}[language=yaml, caption={Limit range example}]
        apiVersion: v1
        kind: ResourceQuota
        metadata:
            name: my-resource-quota
        spec:
            hard:
                request.cpu: 4
                request.memory: 4Gi
                limits.cpu: 10
                limits.memory: 10Gi
    \end{lstlisting}


    \section{Service Account}

    The concept of service accounts is related to other security-related concepts in Kubernetes such as authenatication authorization and RBAC etc.

    There are two types of account in Kubernetes a user account and a service account.
    The user account is used by humans and the service accounts are used by machines.

    A user account could be for an administrator accessing the cluster to performa adminstrictive tasks, or a developer access the cluster to deploy applications etc.

    A service account could be an account used by an appliaiotn to interact with the Kubernetes cluster.
    E.g a monitoring application like Prometheus uses a service account to pull eht Kubernetes API for performance metrics.
    An automated build tools like Jenkins uses service account to deploy applications oin the Kubernetes cluster.

    \begin{lstlisting}[language=yaml, caption={Commands}]
        kubectl create serviceaccount <NAME> # use the create a service account on the CMD
        kubectl get serviceaccount # view service accounts

        kubectl describe serviceaccount <NAME>
    \end{lstlisting}

    When the service account is created it also creates a token automaticlaly.
    The service aaccount token is what must be used bby the external applications while authenaticcating to the Kubernetes API
    The token however is stored as a secret object.

    When a service account is created it first creates the service account objects, then it generateds a token for the service account.
    It then creates a secret object and stores that token inside the secret objects.
    The secret object is then linkedin to the service account.


    To view the token, view the secret oibject by running...

    \begin{lstlisting}[language=yaml, caption={Commands}]
        kubectlr describe serviceaccount <NAME>
    \end{lstlisting}

    This token can then be used as an authentication bearer token while aking a REST call to the Kubernetest API

    You can create a service account and assigne the write premission using role-based access contorl machanisms and export your service accout tokens
    and use it to configure your thrid party application to authetniucat to the kubernetes API.

    Consider if the third-party applicaitn is hosted on the Kuibernetest cluster itself?
    In that case, this whole process of exporting the service account token and configuring the thrid party application to use it can be made simple by automatically moutning the service topken secret as a volume insde the pod hosting the thrid-party application m.
    That way the token to acces the Kubernetes API is already placed inside the pod and can be easily read be the application.
    yo don't have to probvide it manuyally.

    For every namespace in Kubernetes, a service accoutn named default is automatically created.
    Each namespace has its own default service account.
    Whenever a pod is created the default service accounta nd its token are automatically mount to that pod as a volume mount.

    The default service account has restricutions.
    It only has the permission to run basic Kubernetest API queries.
    To use a different service account, modify the pod defnition file to include a service accoutn field and specity the name of the new service account.
    Note: You cannot edit the service account of an exdisting pod. You must delete and recreate the pod.


    Incase of deployment, you will be able to edit the service account as any changes to the pod definition file will automatically trigger a new roll out for the deployment.
    The deployment will take care of deleteing and reacreateing new pods with the right service account.

    Kubernetes automatically mount the default service account if you haven't specified any.
    You may choose not to mount a service account automatically by setting hte automaountServiceaccountToken field to fals int the pod section.

    \subsubsection{Chanages in 1.22 and 1.24 of Kubernetes}
    Every namespace has a default service account.
    That service accoutn has a secret object with a token associate with it.
    When a pod is created, it automatically associates the service account to the pod and mounts the token to a well-known location within the pod.

    This makes the token accessible to a process that's running whithin the pod, and that enables the process to query the Kubernetes API.

    Each JWT requires a separate secret objec tper service account, which results in scalability issues.
    In version 1.22 the TokenRequestAPI was introduced as part of the Kubernetes Enhancement Proposal 12005 that aimed to introduce a mechanism for pvosiioning Kubernetest service account tokens that are more secure and scalable via n aAPI
    Token generated by the token request PAI are audience boudn, they're time bound and object bound.

    \paragraph{Changes in v1.24}
    With version 1.24, a change was made where when you create a service account it no longer automatically creates a secret or a token access sercet.
    You must run the command kubectl create token follwoed by the name ot he service account to generate a token for that service account if you needed one.

    NOTE: see kubernetes service account secrtes


    \section{Taints and Tolerations}

    Taints and tolerations have nothing to do with security or intrusion on the cluster.
    Taints and toleration are used to set restrictions on what pods cna be scehduled on a node.

    Example:
    Consider a cluster where a node contains dedicated resources.
    We do not want the scehduler to put any pod on that node.
    This can be acheived using a taint.
    Pods have no tolerations by default, so they cannot tolerate any taint.

    One half is tainting the node.
    The other is specifying which pods are tolerant to the taint.

    Taints are set on nodes.
    Tolerations are set on pods.

    \begin{lstlisting}[language=yaml, caption={Commands}]
        kubectl taint nodes node-name key=value:taint-effect # taint effect defines what will happen to the pod if they do not tolerat the taint

        kubectl taint nodes node1 app=myapp:NoSchedule
    \end{lstlisting}

    Three types of taint effect
    \begin{itemize}
        \item No scehdule, which mean the pods will not be schdedules on the node.
        \item Preferred no scehdule, which means that the schdeuler/system will try to avoid placing a pod on the node but that is not guaranteed.
        \item Thirs i no execute, which means that the new pods will not be scheduled on the node and existing pods on the node if any will be evicted if they do not tolerate the taint. These pods may hvae been scheduled on the node before the taint was applied to the node.
    \end{itemize}

    \subsubsection{Tolerations - Pods}

    \begin{lstlisting}[language=yaml, caption={Limit range example}]
        apiVersion: v1
        kind: Pod
        metadata:
            name: myapp-pod
        spec:
            containers:
            - name: nginx-container
              image: nginx
            tolerations:
            - key: "app"
              operator: "Equal"
              value: "blue"
              effect: "NoSchedule"
    \end{lstlisting}

    \begin{lstlisting}[language=terminal, caption={Commands}]
        kubectl taint nodes node1 app = blue : NoScehdule
    \end{lstlisting}

    When the pods are no created or updated with the new tolerations, they are either not schedules on nodes or eveicted from the exiting nodes depending on the effect set.

    \subsubsection{No Execute taint effect}
    If a pod does not have a tolerance it is evicted i.e killed.


    Taints and tolerations do not tell a particular pod to go to a particular node.

    Instead it tell the node to accept pods with certain tolerations.

    If the requirement is to restrict a pod to certain nodes it is achieved through the concept of node affinity.

    \subsubsection{Master nodes}
    Taint is automatically applied on this node when the cluster is setup that prevent any pods from being scehduled on this node.
    This behaviour can be modified if required but it is best practice not to deploy applicatn workloads on a master server.

    \begin{lstlisting}[language=terminal, caption={Commands}]
        kubectl describe node kubemaster | grep taint
    \end{lstlisting}


    \section{Node Selectors Logging}
    %TODO see example
    We can set a limitation on the pods so that they only runb on particular nodes.
    There are two ways to do this.
    The first is using node electores.
    For this we look at the pod defintion file.

    \begin{lstlisting}[language=yaml, caption={Node selector example}]
        apiVersion:
        kind: pod
        metadata:
            name: myapp-pod
        spec:
            containers:
            - name: data-processor
              image: data-processor
            nodeSelector:
                size: Large
    \end{lstlisting}

    Kubernetes nodes which nodes are which via key value pairs which are assigned to the node.
    The scehduler uses the labesl to identify the right node.

    To use laberls in a node elector like this, you must have first lableed your nodes prior to creating this pod.

    \begin{lstlisting}[language=terminal, caption={Commands}]
        kubectl label nodes <node-names> <label-key>=<label-value>

        kubectl label nodes node-1 size=Large
    \end{lstlisting}

    Node slectors have limiteations.
    Consider more complext requirements
    Node selectors do not allow for more complex logic.
    For more complextr logic node affinity and anti-affinity features were introduced.


    \section{Node Affinity}
    The primary purpose of node affinity feature si to ensure that pods are hosted on particular nodes.
    This feature allows us with advanved capabilities to limit pod placement on specifi cnodes.

    \begin{lstlisting}[language=yaml, caption={Node affinity example}]
        apiVersion:
        kind: pod
        metadata:
            name: myapp-pod
        spec:
            containers:
            - name: data-processor
              image: data-processor
            affinity:
                nodeAffinity:
                    requiredDuringSchedulingIgnoredDuringExecution:
                        nodeSelectorTerms:
                        - matchExpressions:
                          - key: Size
                            operator: In #Can also use the noIn operator
                            values: # A list of values
                            - Large
                            - Medium
    \end{lstlisting}

    Check documentation for other operators.

    Consider if a node affinity could not match a node with a given expression.

    \subsubsection{Node Affinity Types}
    The type of node affinit defines the behaviour of the scheduler with respect to node affinity and the stages in the lifecycle of the pode.

    There are current two types o node affinity availabled.

    requiredDuringSchedulingIgnoredDuringExecutions
    preferredDuringSchedulingIgnoredDuringExecution

    There are two state in the life cycle of a pod when considering nod affinity during scheduling and during execution.

    Consider that the nodes with matching labels no availble?
    This is where they type of node affinity available comes into play.

    If you slect the required type, which is the first one the scheduler will mandate that the pod be placed on a node with the given affinity reuls.
    If it cannot find on the pod will no be scheduled.
    This type will bes ued in casess where the placement of the pod is crrucial.
    If a matrching node does no exist the pod will no be scheduled.
\hline
    What if the pod placement is less important then running the workload itself?
    In that case you could set it to preferred.
    And in cases where a matching node is not found the scheduler will simply ignore ignored node affinity rules and place the pod on any available node.

    Effectively the scheduler will try bes to place the pod on a matching node but if non are available place it anywhere.
\hline
    The second part of the property or the other state is during execution.
    During execution is the state where a pod has been running and a change is made in the environment that affers node affinity such as a change in the label of a node.

    In the case what would happen to the pods is that value is set to ignored.
    I.e this means pods will continue to run and any changes in the node affinity will not impact them once they are scheduled.

    \section{Taints and Tolerations vs Node Affinity}
    A combination of taints and tolerations an node affinityu rules can be used together to completely dediecate nodes for specifict pods.

    \chapter{Multi-container pods}


    \section{Multi-container pods}
    Idea, decoupling a large monolithic application into microservices allow independently deployable components.
    This architecture helps us scale up and down as well as modify each service as required rather than modifying the entire application.

    At times yo may need two services to work together, such as web server with the main app.

    You don't want to merge and bloat the code of the two services, as each of them target different features and you stil like them to be developed and deployed separately.
    You only need the two to funcitonality to work together.

    You need one web server per main app instance paired together so they can scale up and down together.
    That is why you have multi-container pods that share the same lifecycle, which means they are created and destroyed together.
    They share the same network space, which means they can refere to each other as localhost and they have access to the same storage volumes.

    This way you do not have to eastablish volume sharing or services between the poods to enable communication between them.

    \begin{lstlisting}[language=yaml, caption={Multi container example}]
        apiVersion:
        kind: pod
        metadata:
            name: simple-webapp
            labels:
                name: simple-webapp
        spec:
            containers:
            - name: web-app # This is an array and allows for multiple containers in a single pod
              image: web-app
              ports:
                - containerPort: 8080
            - name: main-app
              image: main-app
    \end{lstlisting}

    \section{Multi-container pods design patterns}
    There are different patterns for multi-container pods.
    \subsubsection{Co-located Containers}
%       TODO see diagram Co-located containers
    Co-located containers - containers located in a single pod. Both the containers are meant to continue to run through the entire pod lifecycle.
    These are usually used when two services are dependent on each other.
    \subsubsection{Init containers}
%    TODO see diagram
    This is used when there are initialization step to be performaed when a pod starts before the applicaiton itself.
    E.g: This could be an init container that waits for a database to be ready before starting the main applicaiton itself.
    The init container does its job and ends its job and then the main applicaiton starts

    \subsubsection{Sidecare container is set up}
    A sidecar container is setup like an init container where the sidecar start first and does its job.
    Instead of ending its run, it continues to run throughout the life cycle of the pod.
    The main application starts after the sidecar container starts.

    This is useful when you have a logger which needs to run along with the main application that needs to start before the main app starts.
    but needs to continue to run as long as the main app runs and then after main app ends.

    NOTE: with co-located containers there is no order as to which container starts. There is no guarantee that one will start without the other.
    If thats not a real rquirement you may use that option.

    NOTE: sidecar containers provides the ability to specify an order of startup and then continue to run throughout the pod lifecycle.

    \begin{lstlisting}[language=yaml, caption={Multi container example}]
        apiVersion:
        kind: Pod
        metadata:
            name: simple-webapp
            labels:
                name: simple-webapp
        spec:
            containers:
            - name: web-app
              image: web-app
              ports:
                - containerPort: 8080
            - name: main-app
              image: main-app
    \end{lstlisting}

    \begin{lstlisting}[language=yaml, caption={Init container example}]
        apiVersion:
        kind: Pod
        metadata:
            name: simple-webapp
            labels:
                name: simple-webapp
        spec:
            containers:
            - name: web-app # This is an array and allows for multiple containers in a single pod
              image: web-app
              ports:
                - containerPort: 8080
            initContainers: # same level as the containers in YAML - initContainers run in sequence
            - name: db-checker #Runs first
              image: busybox
              command: 'wait-for-db-to-start.sh'

            - name: api-checker # Runs second
              image: busybox
              command: 'wait-for-another-api.sh'
    \end{lstlisting}

    \begin{lstlisting}[language=yaml, caption={Init container example}]
        apiVersion: v1
        kind: Pod
        metadata:
            name: myapp-pod
            labels:
                app: myapp
        spec:
            containers:
            - name: myapp-container
              image: busybox:1.28
              command: ['sh', '-c', 'echo The app is running! && sleep 3600']
        initContainers:
            - name: init-myservice
              image: busybox
              command: ['sh', '-c', 'git clone <some-repository-that-will-be-used-by-application> ;']
    \end{lstlisting}

%    See docs about init containers.

    \begin{lstlisting}[language=yaml, caption={Sidecar container example}]
        apiVersion:
        kind: pod
        metadata:
            name: simple-webapp
            labels:
                name: simple-webapp
        spec:
            containers:
            - name: web-app # This is an array and allows for multiple containers in a single pod
              image: web-app
              ports:
                - containerPort: 8080
            initContainers:
            - name: log-shipper
              image: busybox
              command: 'setup-log-shipper.sh'
              retartPolicy: Always
    \end{lstlisting}

%    TODO add Kibana side car example


    \section{Init containers}

    %TODO see docs


    \chapter{Observability}


    \section{Readiness and Liveness Probes}
    \subsubsection{Pod lifecycle}
    Pod status tells us where the pod is in its lifecycle.
    When a pod is first created it is in a pending state.
    This is when the scheduler tries to figure out where to place the pod.
    If the scheduler cannot find the node to place the pod, it remains in a pending state.

    To find out why it's stuck in a pending state run the "kubectl describe pod" command and it will tell you why

    Once the pod is scheduled it goes into a container creating stateus where the images required for the application are pulled and the container starts.

    Once all the containers in a pod start, it goes into a running state, where it continues to be until the program completes successfully or is terminated.

    You can see the pod status in the output of the kubectl get pods command.
    At any point in time the pod status can be one of three values. Pending, Running, ContainerCreating.

    Condition complement pod status.
    It is an array of true or fals values that tell us the state of a pod.

    When a pod is scheduled on a node the pod scheduled condition is set to true.
    WHen the pod is initialised its value is set to true.

    When all the containers in the pod area ready the containers ready condition is set to true and the pod is considered ready itself.


    To see the state of pod conditions run the 'kubectl describe pod' command and look for the conditions section.

    You can also see the ready state of the pod in the output of the kubectl get pods command.

%    --------------------
    The ready conditions indicate the pod is running and is ready to accept user traffic.
    What does that really mean?
    The containers could running different kinds of applicaiton in them.
    It could be a simple script that performs the job or it could be database server or a large web server.
    The script may take a few milliseconds to get ready.

    Scripts may take a few milliseconds to get ready, the DB service make take a few seconds to power up.
    Some webserver could take several minutes to warm up etc.

    By default kuberntes assumes that as soons as the container iscreated it is ready to serve user traffic.
    So it sets the value after the ready condition for each container to true.

    If the applicaiton wihin the container took longer to get ready and the service is unaware of it this cause traffic to hit a pod which is not runnig a live applicaiton.


    We need a way to tie the ready condition to the actual state of the applicaiton inside the container.

    There are different ways that you can define if an applicaitn inside a container is actually ready.
    You can set up different kinds of tests or probes.

    Examples:
    - It could be a HTTP test to see if an API server responds
    - it could be to test if a TCP socket is listening  in the case of a database
    - You could simply execute a command within the container to run a customer script that would exit successfully if the application is ready.
    
    \subsubsection{How to configure a readiness probe}
    \begin{lstlisting}[language=yaml, caption={Sidecar container example}]
        apiVersion:
        kind: Pod
        metadata:
            name: simple-webapp
            labels:
                name: simple-webapp
        spec:
            containers:
            - name: simple-app # This is an array and allows for multiple containers in a single pod
              image: simple-app
              ports:
                - containerPort: 8080

            readinessProbe:
                httpGet:
                    path: /api/ready
                    port: 8080

            initialDelaySeconds: 10
            periodSeconds: 5
            failureThreshold: 8

    \end{lstlisting}

    In the above, kubernetes does not immediately set the ready condition on the container to true.
    Instead, it performs a test to see if the API responds positively.
    Until then, the service does not forward any traffic to the pod as it sees that the pod is not ready.

    There are different ways a probe can be configured.
    For HTTP, use the HTTP get option with the path and the port
    For TCP, use the TCP socket option with the port
    And for executing a command specify the exec option with the command options in an array format.

    There are additional options as well.
    If you know that you rapplication will take a minimum of say 10 seconds to warm up, you can add an additional delay to the probe.
    If you'd like to specify how to probe you can do that using the period seconds options.

    By default, if the application is not raedy after three attempts the probe will stop.
    If you would like more attempts, use the failure threshold option.

    How are readiness probes useful in a multi-pod setup.

    Consider adding a new pod.
    - If the readiness probe has not been set.
    - The service will immediately start routing traffic to the new pod.
    ---- That will result in service disruption to at least some of the users.

    Instead if the pods were configured with the correct readiness probe the service will continue to serve traffic only to the older pods and wait until the new pod is ready.
    Once ready, the traffic will be routed to the new pods as well, ensuring no users are affected.

    \section{Liveness Probes}
    Example:
    Consider a running image of NGINX using Docker and it is serving users.
    For some reason the web server crashes and the NGINX process exists.
    The container exits as well.
    You can see the status of the container with Docker ps command.

    Since dDocker is not an orchestration engine the container conitnues to staty dead and deny services to users until you manually createqa new container.

    This is where Kubernetes orchestration comes in.
    You run the same web applicaiton with Kubernetes.
    Every time the application with Kubernetes crashes, Kubernetes makes an attempt to restart the container to restore service to the users.
    You can see the count of restarts increase in the output of "kube contorl get pods" command


    Consider if the applicaiton no really working but the container continues to stay alive.
    i.e a big in the code or the application is stuck in an infinite loop.
    - As far as kubernetes is concerned the container is up so tha pplicaiton is assumed to be up but the users hitting the container are not served.

    In that case the container needs to be restarted or destroyed and a new container is to be brought up.

    This is where the liveness probe can help us.
    A liveness probe can be configured ont he container to periodicaluy test wheter the applicaiton within the container is actually healthy.
    If the test failes, the container is considered unhealthy and is destroyed and recreated.

    As a developer, you define what it means for an application to be healthy.
    I.e an API is up and running, data base is listening etc. Or a command that performs a test.

    The liveness probe is confiured in the pod defintino ffile. Similar to readiness probe.

    \begin{lstlisting}[language=yaml, caption={Sidecar container example}]
        apiVersion:
        kind: Pod
        metadata:
            name: simple-webapp
            labels:
                name: simple-webapp
        spec:
            containers:
            - name: simple-webapp # This is an array and allows for multiple containers in a single pod
              image: simple-wbeapp
              ports:
                - containerPort: 8080
            livenessProbe:
                httpGet:
                    path: /api/healthy
                    port: 8080

    \end{lstlisting}

%    TODO generate further examples
    Similar to the readiness probe.
    There are HTTP get options for APIS
    TCP sockets for ports
    and exec for commands as well as additional options, like initial delay before the test is run, periodSeconds to define the frequency and success and failure thresholds.

    \section{Logging}

    \begin{lstlisting}[language=terminal, caption={Commands}]
        docker run kodekoud/event-simulator
        docker run -d kodekloud/event-simulator # run in background
        docker logs -f ecf # leaves a live log trail
    \end{lstlisting}

    \begin{lstlisting}[language=yaml, caption={Sidecar container example}]
        apiVersion:
        kind: Pod
        metadata:
            name: event-similator-pod
        spec:
            containers:
            - name: event-simulator
              image: kodekloud/event-simulator
            - name: image-processor
              image: some-image-processor
    \end{lstlisting}

    \begin{lstlisting}[language=terminal, caption={Commands}]
        kubectl create -f event-simulator.yaml
        kubectl logs -f event-simulator.yaml
        kubectl logs -f event-simulator-pod event-simulator # For multiple containers specify the name
    \end{lstlisting}

    These logs are specific to the container running inside the pod.
    A pod however can have multiple containers in it.

    If there are multiple containers within a pod you must specify the name of the container explicitly in the command.
    Otherwise it will fail and ask you to specify a name.

    \section{Monitor and Debug application}
    Kubernetest does not come with a full featured build in monitoring solution (double check as this could be out of date)
    There are a number of open source solutions available today, such as Metrics Server and Prometheus, Elastic stack and properetary solutions like Datadog and Dynatrace.

    Heapster was on of the original project that enabled monitoring and analysis features for Kuberntes.
    This will be seen a lot when looking at reference architectures on monitorign on Kubernetes.
    Heapster is deprecated and a slimmed down version was formed known as the Metrics Server.

    You can have one metric server per Kubernetes cluster.
    The metric server retrieves metrics from each fo the Kubernets nodes and pods, aggregates them and stores them in memeory.

    Note that the metric server is only an in-memoryh montiroing solution and does not store the metrics on the disks.
    As a result you don't see historical performance data.
    For that you must rely on more advanved monitoring solutions.

    \subsubsection{How are metrics generated for pods}
    Kubernetes runs an agent on each node known as Kubelet, which is responsible for receving instricution from the Kubernetest API master server and running pods on the nodes.

    The kublet also contains a subcompnent known as the cAdvisor or Container Advisor.
    cAdvisor is responsible for retreiving performance metrics from pods and exposing them through kublet API to make the metrics available for the metrics server.

    If you're using minikube for your local cluster, run the

    \begin{lstlisting}[language=terminal, caption={Commands}]
        minikube addons enable metrics-server # for local clusters
        # for others
        git clone https://github.com/kubernetes-sigs/metrics-server.git
        kubectl create -f deploy/1.8+/
    \end{lstlisting}

    This command deploys a set of pods services and roles to enable the metric server to poll for performance metrics from the nodes in the cluster.
    Once deployed, give the metric server some time to collect and process data.
    Once process, cluster performance can be viewed by running the command "kubectl top node"
    This proces the CPU and memeory consumption of each of the nodes.
    Can use kubectl top pod to see the performance metrics for the pod. i.e CPU usage and Memory usage

    \chapter{Pod Design}


    \section{Labels, Selectors and Annotations}
    Over time you may end up having hundreds of thousands of Kubernetes objects in your cluster.
    You will need a way to filter and view different object sby different categories, such as to group objects by their type or view objects by their application or by their fuynction.

    You can group and slecto objects using labels and slectors ofr each objects.
    Attach labesl as per needs like app, function etc.

    Then while selecting, specify a condition to filter specific objects.

    \begin{lstlisting}[language=yaml, caption={Selector example}]
        apiVersion:
        kind: Pod
        metadata:
            name: simple-webapp
            labels:
                app: App1
                Function: front-end
        spec:
            containers:
            - name: simple-webapp
              image: simple-webapp
              ports:
                - containerPort: 8080
    \end{lstlisting}

    \begin{lstlisting}[language=terminal, caption={Commands}]
        kubectl get pods --selector app=App1
    \end{lstlisting}

    This is one use case of labels and selectors.
    Kubenertes use labels and slectors internally to connect different objects together.

    \begin{lstlisting}[language=yaml, caption={ReplicaSet labels example}]
        apiVersion:
        kind: Pod
        metadata:
            name: simple-webapp
            labels: # replica set labels
                app: App1
                Function: front-end
        spec:
            replicas: 3
            selector:
                matchLabels:
                    app: App1 # used to connect replicaset to the pods
            template:
                metadata:
                    labels: # pods labels
                        app: App1
                        function: Front-end
            spec:
                containers:
                - name: simple-webapp
                  image: simple-webapp
    \end{lstlisting}

    In order to connect the replica set to the pods, we configure the selector field under Replica set specification to match the labels defined on the pod.
    A single label will do if it matches correctly.

    If there are other pod poarts with same label but with a different function, then you coudl specify both the labels to ensure that the right parts are discovered by the replica set on createion.
    If the labels match, the replica set is created successfully.
    
    It works the same for other objects like a service.
    When a service is created, it uses the selector defined in the service definitino file to match the labels set on the pods in the replica set defintion file.
    
    \subsubsection{Annotations}
    While labels an selectors are used to group and select objects, annotations are used to record other details for informatory purposes.

    \begin{lstlisting}[language=yaml, caption={ReplicaSet labels example}]
        apiVersion: apps/v1
        kind: ReplicaSet
        metadata:
            name: simple-webapp
            labels:
                app: App1
                function: Front-end
        annoations:
            buildversions: 1.34 # These are details that can just be used for some other purpose
        spec:
            replicas: 3
            selector:
                matchLabels:
                    app: App1
            template:
                metadata:
                    labels:
                        app: App1
                        function: Front-end
                spec:
                    containers:
                    - name: simple-webapp
                      image: simple-webapp
    \end{lstlisting}

    \section{Rolling updates and Rollbacks in Deployments}
    When you first create a deployment it triggeres a rollout.
    A new rollour creates a new replica set, which is recorded as a new deployment revision in the future when the applicaiton is upgraded.
    This means when the container version is updated to a new one, a new rollout is triggered and new deployrevision is created named revision two.

    This helps keep track fot he changess made to out deployment and enabled us to roll back to a previous version fo deployment if necessary.

    This comand can be used to see the status of a rollout.
    \begin{lstlisting}[language=terminal, caption={Commands}]
        kubectl rollout status deployment/myapp-deployment
        kubectl rollout history deployment/myapp-deployment # To see the revions and history of a rollout, run this command.
    \end{lstlisting}

    \subsubsection{Deployment strategies}
    There are two types of deployment strategies.

    Examples: say you have give replicas of your web application instance deployed.
    One wasy to updated these to a newer verions is to destroy all of these and then create newer versions of application instances.
    I.e destroy the five instances and then deploy five new version of the instances.

    The problem here is that during the gap where they are destroyed and the new versions are up, the application is down and inaccessible to users.

    This strategy is know as the recreate strategy and is not the default deployment strategy.

    The second strategy is where we don not destroy the all at once.
    Instead we take down the older verions and bring up a newer version one by one.
    This way the application never goes down and the upgrade is seamless.

    Note: If you do not specify a strategy while create the deployment, it will assume it to be a rolling update.
    In other words, the rolling update is the default deployment strategy.

    How exactyla do yo uupdate your deploiyment?
    This could be different thing, such as updating your applicaiton verions by update the verions of Docker containers used, updating their labels or updating the number of replicas etc.
    Since there is a deployment definiton file, it is easy for use to modify this file once we make the necessary changes.

    \begin{lstlisting}[language=yaml, caption={Deployment example}]
        apiVersion: apps/v1
        kind: Deployment
        metadata:
            name: myapp-deployment
            labels:
                app: myapp
                function: front-end
        spec:
            template:
                metadata:
                    name: myapp-pod
                    labels:
                        app: myapp
                        type: front-end
                spec:
                    containers:
                    - name: nginx-container
                      image: nginx: 1.7.1
        replicas: 3
        selector:
            matchLabels:
                type: front-end
    \end{lstlisting}

    \begin{lstlisting}[language=terminal, caption={Commands}]
        kubectl apply -f deployment-definition
    \end{lstlisting}

    Run the above command to apply the changes.
    A new rollout is triggered and a new revision of the deployment is created.

    There is another way to do the same thing.
    You could use the kube conttorl set image command to update the image of your application.

    \begin{lstlisting}[language=terminal, caption={Commands}]
        kubectl set image deployment/myapp-deployment \nginx=nginx:1.9.1
    \end{lstlisting}

    Remeber, doing it this way will result in the deployment definiton file having a different configuraiton.
    You msut be carefuy when using the same defition file to make changes in the future.

    The difference betwen the createate and rolling update strategies can also be seen when you view the deployments in details from kube control.

    \begin{lstlisting}[language=terminal, caption={Commands}]
        kubectl describe deployment myapp-deployment
    \end{lstlisting}

    Whens using the above command yyou will notice when the recarete strategy was used.
    The events indivate that the old replica set was scaled down to zero first and then the new replica sets scaled up to give.

    However, when the rolling update strategy was used, the old replica set was scaled down one at a time,
    simulataneously scaling up the new replica set one at a time.

    \subsubsection{Upgrades}
    Example: deploy five replicas.
    First creates a replica set automatically, which int turn creates the number of pods required to meet the number of replicas.
    When you upgrade your application, the Kubernetes deployyment object creates a new replica set under the fhhood and starts deploying the containerrs there.
    At the same time taking down the pods in the old replica set following a rolling update strategy.

    This can be seen when you try to list the replica sets using "kubetl get replicasets" command.

    This will show the old replica set with zero pods and the new replica set with five parts.

    Scenario:
    Say for instance, once you upgrade your application you realise something isn't right. Something is wrong with the new version of the build you used etc.
    You now want to roll back your update.

    Kubernetes deployments allow you to roll back to a previous revision, to undoi a change.


    Run the following command with the name of the deployment
    \begin{lstlisting}[language=terminal, caption={Commands}]
        kubectl rollout undo deployment myapp-deployment
    \end{lstlisting}

    The deployment will then destroy the pods in the new replica set and bring the olders ones up in the old replica set and you application is back to its older format.

    When you compare the output of the kubectl get replicasets command before and after the rollback. You will notice the difference between before the rollback.

    \subsection{Command summary}
    \begin{lstlisting}[language=terminal, caption={Sumnmary of commands}]
        kubectl create -f deployment-defintion.yml # Create
        kubectl get deployments # Get

        # ---- update
        kubectl apply -f deployment-definition.yml
        kubectl set image deployment/myapp-deployment nginx=nginx:1.9.1

        # ----status
        kubectl rollout status deployment/myapp-deployment
        kubectl rollout history deployment/myapp-deployment

        kubectl rollout undo deployment/myapp-deployment # Rollback!

    \end{lstlisting}


    \section{Updating a deployment}
%    TODO see examples

    \section{Deployment strategy - Blue Green}
    There are other strategies as well that can be specifed in the deployment.
    One of them is blue-green deployments.

    Here we have the new version deployed alongside the old version.
    The old version is called blue and the new version is called green.
    100\% of the traffic is still routed to the old veriosns.
    At this point in time tests are run on the new version and once all the tests have passed we switch traffic to the new version all at once.

    These strategies are best implemented with service measures like Istio.
    See service mesh with istio.

    Heres how it works.
    First we have the originoal verison of our applicaiton deployed as a service.
    We call it the blue deployment.
    Then we create a service to route traffic to it.

    Now to associate ther service to the pods in the deployment we set a label on the pods. Lets call it version set to v1.
    We use the same label as the slector on ther service.

    We then deploy a second deployment.
    Lets call it green with a new version of the application.
    Once all tests are passed we route traffic from the service to the green deployment by wsiwtching the label selceotr on the service.
    So for the new deployment, we set a label version to v2.
    Now that we have set the label version to v2 to the pods on the new deployment all we need to do is switch the labels specified under the slector of the service to v2 as well.
    Then the service switches traffic to the pods in the green deployment.

    \subsection{Examples}

    \begin{lstlisting}[language=yaml, caption={Blue deployment example}]
        apiVersion: apps/v1
        kind: Deployment
        metadata:
            name: myapp-blue
            labels:
                app: myapp
                type: front-end
        spec:
            template:
                metadata:
                    name: myapp-pod
                    labels:
                        versions: v1 # Important for the deployment
                spec:
                    containers:
                    - name: app-container
                      image: myapp-image:1.0
        replicas: 5
        selectors:
            matchLabels:
                version: v1
    \end{lstlisting}

    \begin{lstlisting}[language=yaml, caption={Green deployment example}]
        apiVersion: apps/v1
        kind: Deployment
        metadata:
            name: myapp-green
            labels:
                app: myapp
                type: front-end
        spec:
            template:
                metadata:
                    name: myapp-pod
                    labels:
                        versions: v2 # Important for the deployment
                spec:
                    containers:
                    - name: app-container
                      image: myapp-image:2.0
        replicas: 5
        selectors:
            matchLabels:
                version: v2
    \end{lstlisting}

    \begin{lstlisting}[language=yaml, caption={Service defintion example}]
        apiVersion: apps/v1
        kind: Service
        metadata:
            name: my-service
        spec:
            selector:
                version: v1 #Change this to v2 when ready to reroute traffic to the green deployment
    \end{lstlisting}

    \section{Deployment strategy - Canary}
    In this strategy we deploy the new version and route only a small percentrage of traffic to it.
    The majority of traffic is being routed to the older version, but we have a small percentage routed to the new version.
    At this point, we run tests and if everything looks good we upgrade the origonal deployment with the newer version of the application.
%TODO see/create diagram

    With the canary deployment we want traffic to go to both deployments at the same time.
    However we only want to route a small percentage of traffic to version two, so how do we do that?

    For this we can create a common label.
    Call it App, which is set to frontend and we update the slector label in the service to match this common label.
    That solves the first problem.

    We now have traffic going to both the vewions of applicaitons but it's reouting it equally. I.e 50\% to each deployment.

    How do we reduce traffic to canary deployemnt?
    We only want a small percentage of traffic to go to the canary deployment.

    We can do this by simply reducein the number of pods in the canary deployment, to the minimum possible, in this case one.
    So since a service distirbutes traffic between all pods equally the canary deployment will recieve less traffic.

    Now once tests are prefomred and we're confdent with the new version, we can now upgrade the version of pods in the primary deployment and delete the canary deployment altogether.

    one of the caveats of performing a canary update this way with just Kubernetes deployments and service is that we have limited contorl over the split of traffic between each deployment.
    The traffic split is always going to be governed by the number of pods present in each deployment.

    Notes about Istio:
    With Istio, you can define the exact percentage of traffic to be distributed between each deployment
    and its not really dependent upon the number of pods in a deployment.

    \subsection{Canary Deployment Examples}

    \begin{lstlisting}[language=yaml, caption={Primary deployment example}]
        apiVersion: apps/v1
        kind: Deployment
        metadata:
            name: myapp-primary
            labels:
                app: myapp
                type: front-end
        spec:
            template:
                metadata:
                    name: myapp-pod
                    labels:
                        version: v1 # Important for the deployment
                        app: front-end
                spec:
                    containers:
                    - name: app-container
                      image: myapp-image:1.0
        replicas: 5
        selectors:
            matchLabels:
                app: front-end
    \end{lstlisting}

    \begin{lstlisting}[language=yaml, caption={Canary deployment example}]
        apiVersion: apps/v1
        kind: Deployment
        metadata:
            name: myapp-canary
            labels:
                app: myapp
                type: front-end
        spec:
            template:
                metadata:
                    name: myapp-pod
                    labels:
                        versions: v2 # Important for the deployment
                        app: front-end
                spec:
                    containers:
                    - name: app-container
                      image: myapp-image:2.0
        replicas: 1 # Set to one since we only want a small percentage of traffic
        selectors:
            matchLabels:
                version: front-end
    \end{lstlisting}

    \begin{lstlisting}[language=yaml, caption={Service defintion example}]
        apiVersion: v1
        kind: Service
        metadata:
            name: my-service
        spec:
            selector:
                app: front-end
    \end{lstlisting}


    \section{Jobs}
    There are different types of workloads that a container can servce. E.g web, applicaiton and database.
    Some workloads are meant to conitue to run for a long period of time until manually taken down.

    There are other kinds of workloads as match processing, analytics or reporting which carry out a specific task and then finish.
    Ee.g performing a computation or processing an image, performing some kind analytics on a large data set, generating a report and sending an email etc.

    These workloads are meant to live for a short period of time, perform a set of tasks and then finish.


    Consider in docker.

    Docker run ubuntu expr 3 + 2.
    Docker runs the expression, produced an output and the container then goes into a exited state.
    This can be seen with docker ps -a. Return code of the operation is shown in the bracket as well.

    \begin{lstlisting}[language=yaml, caption={Example}]
        apiVersion: apps/v1
        kind: Deployment
        metadata:
            name: math-pod
        spec:
            containers:
            - name: math-add
              image: ubuntu
              command: ['expr', '3', '+', '2']
    \end{lstlisting}

    \begin{lstlisting}[language=terminal, caption={Sumnmary of commands}]
        kubectl create -f pod-definition.yaml
        kubectl get pods
    \end{lstlisting}

    Replicating the same thing in kubernetes.
    Create a pod defintionfile to perform the same operation.
    When the pod is crated, it runs a container, performs the computation taask and exist.
    The pod goes into a completed state but it then recreates the container in an attempt to leave it running.
    Againt the container peforms the rquired computation task and exists
    This continue to happen un til a threshold is reached.

    Why does this happen?

    Kubernetes wants your application to live forever.
    The default behaviour of pods is to attempt to restart the container in an effort to keep it running.
    This behaviour is defined by the restart policy set on the pod which is by default set to always.
    That is why pod always recreates the container when it exists.

    \begin{lstlisting}[language=yaml, caption={Example}]
        apiVersion: apps/v1
        kind: Deployment
        metadata:
            name: math-pod
        spec:
            containers:
            - name: math-add
              image: ubuntu
              command: ['expr', '3', '+', '2']
            restartPolicy: Always # Can set to 'Never' or 'onFailure'
    \end{lstlisting}

    You can override this behaviour by setting this property to never or on failure.
    That way kibernetes does not retard the container once the job is finished.

    We have new use cases for batch processing.
    We have large datasets that require multiple pods to process the data in parallel.
    - We want to make sure that all pods perform the task assigned ot them successfully and then exit.
    We need a managed that can create as many pods as we want to get work done and snure work gets done successfully.

    \subsection{ReplicaSets vs Jobs}
    While a replica set is used to make sure a specfied number of pods are running at all times.
    A job is used to run a set of pods to perform a given task to completion.

    \subsubsection{Example}
    \begin{lstlisting}[language=yaml, caption={Pod definition file}]
        apiVersion: v1
        kind: Pod
        metadata:
            name: math-pod
        spec:
            containers:
            - name: math-add
              image: ubuntu
              command: ['expr', '3', '+', '2']
            restartPolicy: Never
    \end{lstlisting}

    \begin{lstlisting}[language=yaml, caption={Job definition file}]
        apiVersion: batch/v1
        kind: Job
        metadata:
            name: math-add-job
        spec:
            containers:
            - name: math-add
              image: ubuntu
              command: ['expr', '3', '+', '2']
            restartPolicy: Never
    \end{lstlisting}

    Use these commands to seet that kubernetes did not try and restart the pod.

    \begin{lstlisting}[language=terminal, caption={Sumnmary of commands}]
        kubectl create -f job-definition.yaml
        kubectl get jobs
        kubectl get pods
    \end{lstlisting}

    Output of the pod should be in the pod's standard output.
    The standard output of a container can be seen using hte logs command.....

    \begin{lstlisting}[language=terminal, caption={Sumnmary of commands}]
        kubectl logs math-add-job-1d87pn # See the output
        kubectl delete job math-add-job # Delete the job
    \end{lstlisting}

    Deleting the job will also result in deleting the pods that were created by the job.

    Note: Typically this is not how jobs are implemented in the read world.

    \subsection{Running multiple pods}

    \begin{lstlisting}[language=yaml, caption={Job definition file}]
        apiVersion: batch/v1
        kind: Job
        metadata:
            name: math-add-job
        spec:
            completions: 3
            template:
                spec:
                    containers:
                    - name: math-add
                      image: ubuntu
                      command: ['expr', '3', '+', '2']
                    restartPolicy: Never
    \end{lstlisting}

    \begin{lstlisting}[language=terminal, caption={Sumnmary of commands}]
        kubectl create -f job-defintion.yaml
        kubectl get jobs
        kubectl get pods
    \end{lstlisting}

    We should see the desired count is three and the successful count is three.
    By default these pods are created on after another.
    The second part is created only after the first is finished.

    \subsubsection{Failing jobs}
    Consider if the pods are going to fail?
%    TODO pinch random container example
    \begin{lstlisting}[language=yaml, caption={Job definition file}]
        apiVersion: batch/v1
        kind: Job
        metadata:
            name: random-error-job
        spec:
            completions: 3
            template:
                spec:
                    containers:
                    - name: random-error
                      image: kodekloud/random-error
                    restartPolicy: Never
    \end{lstlisting}

\subsubsection{Creating jobs in parallel}
    Jobs can be created in paralellism
    \begin{lstlisting}[language=yaml, caption={Job definition file}]
        apiVersion: batch/v1
        kind: Job
        metadata:
            name: random-error-job
        spec:
            completions: 3
            parallelism: 3
            template:
                spec:
                    containers:
                    - name: random-error
                      image: kodekloud/random-error
                    restartPolicy: Never
    \end{lstlisting}

    \begin{lstlisting}[language=terminal, caption={Sumnmary of commands}]
        kubectl create -f job-defintion.yaml
        kubectl get jobs
        kubectl get pods
    \end{lstlisting}

%    TODO see example

    \section{CronJobs}
    A CronJob is a job that can be scheduled just like CronTab in Linux.

    Example: say you have a job that generates a report and sends an email. You can create the job using the cube control create command but it runs instantly .

    Instead, you could create a CronJob to schedule and run it periodically.

    \subsection{CronJob Example}

    \begin{lstlisting}[language=yaml, caption={Job definition file}]
        apiVersion: batch/v1beta1
        kind: CronJob
        metadata:
            name: reporting-cron-job
        spec:
            schedule: "*/1 * * * *"
            jobTemplate:
                spec:
                    completions: 3
                    parallelism: 3
                    template:
                        spec:
                            containers:
                            - name: reporting-tool
                              image: reporting-tool
                            retartPolicy: Never
    \end{lstlisting}

    \begin{lstlisting}[language=terminal, caption={Sumnmary of commands}]
        kubectl create -f cron-job-definition.yaml
        kubectl get cronjob
    \end{lstlisting}


    \chapter{Services and Networking}

    \section{Network Policies}

    \section{Developing network policies}

    \section{Ingress networking}

    \chapter{State Persistence}

    \section{Introduction to Docker Storage}

    \section{Storage in Docker}

    \section{Volume Driver Plugins in Docker}

    \section{Volumes in Kubernetes}

    \section{Persistent Volumes}

    \section{Persistent volume claims}

    \section{Using PVCs in Pods}

    \section{Storage Classes}

    \section{Why Stateful Sets}

    \section{Sateful sets introduction}

    \section{Headless Services}

    \section{Storage in StatefulSets}

    \chapter{Security}

    \section{Kubernetes Security Primitives}

    \section{Authentication}

    \section{KubeConfig}

    \section{API Groups}

    \section{Authorization}

    \section{Role based access controls}

    \section{Cluster Roles}

    \section{Admission Controllers}

    \section{Validating and Mutating Admission Controllers}

    \section{API Versions}

    \section{API Deprecations}

    \section{Custom Resource Definition}

    \section{Custom Controllers}

    \section{Operator Framework}

    \chapter{Helm Fundementals}

    \section{Helm Introduction}

    \section{Install Helm}

    \section{Helm concepts}

    \chapter{Kustomize}

    \section{Kustomize Problem Statement and Ideology}

    \section{Kustomize vs Helm}

    \section{Kustomization yaml file}

    \section{Kustomization Output}

    \section{Kustomization ApiVersion and Kind}

    \section{Managing Directories}

    \section{Common Transformers}

    \section{Image Transformers}

    \section{Patches Intro}

    \section{Different Types of Patches}

    \section{Patches Dictionary}

    \section{Patches list}

    \section{Overlays}

    \section{Components}

\end{document}